% Options for packages loaded elsewhere
% Options for packages loaded elsewhere
\PassOptionsToPackage{unicode}{hyperref}
\PassOptionsToPackage{hyphens}{url}
\PassOptionsToPackage{dvipsnames,svgnames,x11names}{xcolor}
%
\documentclass[
  8pt,
  letterpaper,
  DIV=11,
  numbers=noendperiod]{scrreprt}
\usepackage{xcolor}
\usepackage[margin=0.75in]{geometry}
\usepackage{amsmath,amssymb}
\setcounter{secnumdepth}{5}
\usepackage{iftex}
\ifPDFTeX
  \usepackage[T1]{fontenc}
  \usepackage[utf8]{inputenc}
  \usepackage{textcomp} % provide euro and other symbols
\else % if luatex or xetex
  \usepackage{unicode-math} % this also loads fontspec
  \defaultfontfeatures{Scale=MatchLowercase}
  \defaultfontfeatures[\rmfamily]{Ligatures=TeX,Scale=1}
\fi
\usepackage{lmodern}
\ifPDFTeX\else
  % xetex/luatex font selection
\fi
% Use upquote if available, for straight quotes in verbatim environments
\IfFileExists{upquote.sty}{\usepackage{upquote}}{}
\IfFileExists{microtype.sty}{% use microtype if available
  \usepackage[]{microtype}
  \UseMicrotypeSet[protrusion]{basicmath} % disable protrusion for tt fonts
}{}
\makeatletter
\@ifundefined{KOMAClassName}{% if non-KOMA class
  \IfFileExists{parskip.sty}{%
    \usepackage{parskip}
  }{% else
    \setlength{\parindent}{0pt}
    \setlength{\parskip}{6pt plus 2pt minus 1pt}}
}{% if KOMA class
  \KOMAoptions{parskip=half}}
\makeatother
% Make \paragraph and \subparagraph free-standing
\makeatletter
\ifx\paragraph\undefined\else
  \let\oldparagraph\paragraph
  \renewcommand{\paragraph}{
    \@ifstar
      \xxxParagraphStar
      \xxxParagraphNoStar
  }
  \newcommand{\xxxParagraphStar}[1]{\oldparagraph*{#1}\mbox{}}
  \newcommand{\xxxParagraphNoStar}[1]{\oldparagraph{#1}\mbox{}}
\fi
\ifx\subparagraph\undefined\else
  \let\oldsubparagraph\subparagraph
  \renewcommand{\subparagraph}{
    \@ifstar
      \xxxSubParagraphStar
      \xxxSubParagraphNoStar
  }
  \newcommand{\xxxSubParagraphStar}[1]{\oldsubparagraph*{#1}\mbox{}}
  \newcommand{\xxxSubParagraphNoStar}[1]{\oldsubparagraph{#1}\mbox{}}
\fi
\makeatother


\usepackage{longtable,booktabs,array}
\usepackage{calc} % for calculating minipage widths
% Correct order of tables after \paragraph or \subparagraph
\usepackage{etoolbox}
\makeatletter
\patchcmd\longtable{\par}{\if@noskipsec\mbox{}\fi\par}{}{}
\makeatother
% Allow footnotes in longtable head/foot
\IfFileExists{footnotehyper.sty}{\usepackage{footnotehyper}}{\usepackage{footnote}}
\makesavenoteenv{longtable}
\usepackage{graphicx}
\makeatletter
\newsavebox\pandoc@box
\newcommand*\pandocbounded[1]{% scales image to fit in text height/width
  \sbox\pandoc@box{#1}%
  \Gscale@div\@tempa{\textheight}{\dimexpr\ht\pandoc@box+\dp\pandoc@box\relax}%
  \Gscale@div\@tempb{\linewidth}{\wd\pandoc@box}%
  \ifdim\@tempb\p@<\@tempa\p@\let\@tempa\@tempb\fi% select the smaller of both
  \ifdim\@tempa\p@<\p@\scalebox{\@tempa}{\usebox\pandoc@box}%
  \else\usebox{\pandoc@box}%
  \fi%
}
% Set default figure placement to htbp
\def\fps@figure{htbp}
\makeatother





\setlength{\emergencystretch}{3em} % prevent overfull lines

\providecommand{\tightlist}{%
  \setlength{\itemsep}{0pt}\setlength{\parskip}{0pt}}



 


\usepackage[section]{placeins}
\usepackage{multirow}
\usepackage{booktabs}
\usepackage{adjustbox}
\KOMAoption{captions}{tableheading}
\usepackage{textcomp}
\makeatletter
\@ifpackageloaded{bookmark}{}{\usepackage{bookmark}}
\makeatother
\makeatletter
\@ifpackageloaded{caption}{}{\usepackage{caption}}
\AtBeginDocument{%
\ifdefined\contentsname
  \renewcommand*\contentsname{Table of contents}
\else
  \newcommand\contentsname{Table of contents}
\fi
\ifdefined\listfigurename
  \renewcommand*\listfigurename{List of Figures}
\else
  \newcommand\listfigurename{List of Figures}
\fi
\ifdefined\listtablename
  \renewcommand*\listtablename{List of Tables}
\else
  \newcommand\listtablename{List of Tables}
\fi
\ifdefined\figurename
  \renewcommand*\figurename{Figure}
\else
  \newcommand\figurename{Figure}
\fi
\ifdefined\tablename
  \renewcommand*\tablename{Table}
\else
  \newcommand\tablename{Table}
\fi
}
\@ifpackageloaded{float}{}{\usepackage{float}}
\floatstyle{ruled}
\@ifundefined{c@chapter}{\newfloat{codelisting}{h}{lop}}{\newfloat{codelisting}{h}{lop}[chapter]}
\floatname{codelisting}{Listing}
\newcommand*\listoflistings{\listof{codelisting}{List of Listings}}
\makeatother
\makeatletter
\makeatother
\makeatletter
\@ifpackageloaded{caption}{}{\usepackage{caption}}
\@ifpackageloaded{subcaption}{}{\usepackage{subcaption}}
\makeatother
\usepackage{bookmark}
\IfFileExists{xurl.sty}{\usepackage{xurl}}{} % add URL line breaks if available
\urlstyle{same}
\hypersetup{
  pdftitle={Supervised Classification for anomaly detection in IoT devices using RT-IoT2022 dataset},
  pdfauthor={Joe Nguyen},
  colorlinks=true,
  linkcolor={blue},
  filecolor={Maroon},
  citecolor={Blue},
  urlcolor={Blue},
  pdfcreator={LaTeX via pandoc}}


\title{Supervised Classification for anomaly detection in IoT devices
using RT-IoT2022 dataset}
\author{Joe Nguyen}
\date{}
\begin{document}
\maketitle

\renewcommand*\contentsname{Table of contents}
{
\hypersetup{linkcolor=}
\setcounter{tocdepth}{2}
\tableofcontents
}

\bookmarksetup{startatroot}

\chapter{Introduction:}\label{introduction}

The rapid adoption of Internet of Things (IoT) technologies across
sectors such as healthcare, manufacturing, and agriculture has increased
reliance on continuous connectivity and large-scale data exchange. While
this connectivity enables new capabilities, it also expands the attack
surface, making IoT devices attractive entry points for cyberattacks.
Compromised devices can be leveraged to infiltrate broader network
infrastructures through vulnerable endpoints (Sharmila and Nagapadma,
2023). The financial impact of these attacks continues to grow. The IBM
Cost of a Data Breach Report 2025 estimates the global average cost of a
breach at USD 4.44 million, highlighting the sustained economic risk
posed by increasingly sophisticated and AI-enabled threats (IBM
Security, 2025). Real-world incidents illustrate these risks clearly. In
2021, the Verkada breach exposed live feeds from approximately 150
million surveillance cameras, and in a separate incident, an attacker
manipulated chemical levels at a Florida water treatment facility.
Together, these events demonstrate the real-world consequences of
cyberattacks targeting IoT systems and reinforce the need for reliable
anomaly detection mechanisms as IoT deployments continue to scale
(Sharmila and Nagapadma, 2023).

This project replicates and extends the work of Sharmila and Nagapadma
(2023) by evaluating autoencoder-based anomaly detection for identifying
malicious IoT network traffic. Beyond replication, the study emphasizes
systematic model selection, rigorous evaluation of predictive and
generalization performance, and interpretation of how learned
representations relate to underlying traffic features. Because IoT
environments exhibit heterogeneous and device-specific traffic patterns,
effective anomaly detection requires models that learn normal behavioral
baselines rather than predefined attack signatures. Consistent with
prior work, autoencoder models are trained exclusively on benign traffic
under the assumption that anomalous activity produces higher
reconstruction error. Experiments are conducted using real-time traffic
collected from four representative devices---ThingSpeak-LED, MQTT-Temp,
Amazon Alexa, and Wipro Bulb---allowing evaluation of model robustness
across diverse IoT operating conditions.

Modern intrusion detection systems (IDS) for IoT environments
increasingly rely on unsupervised learning, particularly anomaly
detection, to identify rare deviations from normal network behavior. As
discussed by Sharmila and Nagapadma (2023), this approach is well suited
to IoT settings where labeled attack data is limited and anomalous
events may have significant operational or environmental consequences.
Progress in this area, however, has been constrained by the scarcity of
realistic IoT network traffic datasets. To address this limitation, the
authors emphasize the importance of datasets derived from real-time IoT
devices that capture both benign and malicious activity patterns.

The authors also note that deploying AI-based IDS frameworks directly on
IoT devices introduces practical constraints due to limited memory,
processing power, and energy availability. Consequently, conventional
deep learning models are often impractical without optimization. Their
work highlights model optimization techniques---including network
pruning and numerical quantization---as effective strategies for
reducing model size and computational overhead while maintaining
detection performance. These optimizations enable faster inference,
lower energy consumption, and more feasible deployment of anomaly
detection models on resource-constrained IoT hardware.

\section{Project Objectives:}\label{project-objectives}

The goal of this project is to evaluate machine learning approaches for
detecting anomalous network behavior in IoT environments using the
RT-IoT 2022 dataset, with emphasis on both predictive performance and
real-world deployment constraints.

The project is guided by the following objectives:

\begin{enumerate}
\def\labelenumi{\arabic{enumi}.}
\item
  Establish a Supervised Baseline

  \begin{itemize}
  \item
    Develop a Logistic Regression classifier as an interpretable
    baseline for distinguishing between normal and attack network
    traffic.
  \item
    This model provides a statistical reference point for evaluating
    more complex approaches and helps assess how effectively a linear
    decision boundary separates benign and malicious IoT activity.
  \end{itemize}
\item
  Develop Autoencoder-Based Anomaly Detection Models

  \begin{itemize}
  \item
    Implement autoencoder architectures (QAE-u8 and QAE-f16) trained
    primarily on normal network traffic, using reconstruction error to
    identify anomalous behavior.
  \item
    These models aim to learn compact representations of normal device
    activity and flag deviations that may correspond to cyberattacks.
    The analysis examines how representation learning compares with
    classical supervised classification for intrusion detection.
  \end{itemize}
\item
  Evaluate Edge Deployment Under Resource Constraints

  \begin{itemize}
  \item
    Deploy trained models on a Raspberry Pi to evaluate feasibility
    within resource-constrained IoT environments.
  \item
    Evaluation focuses on operational characteristics, including:

    \begin{itemize}
    \item
      inference latency
    \item
      CPU and memory utilization
    \item
      power and computational efficiency
    \item
      model size and deployment overhead
    \end{itemize}
  \item
    This objective connects model performance with practical
    edge-computing requirements.
  \end{itemize}
\item
  Analyze Performance--Efficiency Tradeoffs

  \begin{itemize}
  \item
    Compare supervised and unsupervised approaches not only by
    predictive performance but also by efficiency when deployed on
    constrained hardware.
  \item
    The goal is to identify which modeling strategy provides the best
    balance between:

    \begin{itemize}
    \item
      detection capability
    \item
      computational cost
    \item
      real-world deployability in IoT security settings
    \end{itemize}
  \end{itemize}
\end{enumerate}

\section{Data Summary:}\label{data-summary}

Source Link: https://archive.ics.uci.edu/dataset/942/rt-iot2022

RT-IoT2022 is a network traffic dataset collected from a real-time IoT
environment containing both normal device activity and multiple
simulated cyberattacks. The dataset includes traffic generated by IoT
devices alongside attack scenarios such as brute-force SSH attempts,
DDoS attacks (e.g., Hping and Slowloris), and network scanning patterns.
Network flows are captured bidirectionally using the Zeek monitoring
framework with the Flowmeter plugin, producing structured flow-level
features that characterize packet behavior, timing, and protocol
dynamics.

In this project, RT-IoT2022 serves as the foundation for developing and
evaluating intrusion detection approaches---particularly anomaly
detection using autoencoders---aimed at identifying malicious behavior
within real-time IoT network traffic.

\begin{enumerate}
\def\labelenumi{\arabic{enumi})}
\item
  \textbf{Predictor (Feature) Variables}: The RT-IoT2022 dataset
  contains \textbf{83 predictor features} describing network traffic
  flows generated by real-time IoT devices. Rather than representing
  unrelated variables, these features follow a systematic naming
  convention built from combinations of directional prefixes and
  statistical or behavioral suffixes. Each feature captures a measurable
  property of packet flows---such as volume, timing, payload
  characteristics, protocol behavior, or connection state---computed
  either for one traffic direction or for the entire bidirectional flow.

  This structured design allows the dataset to encode network behavior
  at multiple granularities (packet, flow, and temporal statistics),
  making it well suited for intrusion detection and anomaly detection
  tasks. Understanding the prefix--suffix pattern is essential for
  interpreting the predictors.

  \textbf{Prefixes (traffic scope / direction)}:

  \begin{itemize}
  \tightlist
  \item
    fwd\_ (forward): statistics computed for packets traveling from the
    origin/source to the responder/destination.
  \item
    bwd\_ (backward): statistics computed for packets traveling from the
    responder/destination back to the origin/source.
  \item
    flow\_ (flow-level): aggregate statistics computed across the entire
    bidirectional communication flow.
  \item
    active\_ / idle\_: metrics describing active transmission periods
    versus inactive (silent) intervals within a connection.
  \item
    (No prefix): global or contextual attributes describing protocol,
    ports, or overall flow characteristics.
  \end{itemize}

  \textbf{Suffixes (measurement type / statistical summary)}:

  \begin{itemize}
  \tightlist
  \item
    .min --- minimum observed value within the flow.
  \item
    .max --- maximum observed value within the flow.
  \item
    .tot --- total accumulated value.
  \item
    .avg --- arithmetic mean across packets/events.
  \item
    .std --- standard deviation measuring variability.
  \item
    \_per\_sec --- rate-based metric normalized by time.
  \item
    \_flag\_count --- count of packets containing specific TCP control
    flags.
  \item
    \_bytes, \_pkts, \_payload --- volume-based measurements describing
    byte counts, packet counts, or payload sizes.
  \item
    \_iat --- inter-arrival time statistics between packets.
  \item
    \_window\_size --- TCP window characteristics reflecting
    transport-layer behavior.
  \end{itemize}

  This prefix--suffix structure enables consistent interpretation across
  the 83 features while allowing the same behavioral concept (e.g.,
  payload size or timing variability) to be analyzed across traffic
  directions and statistical summaries.
\item
  \textbf{Outcome (Response) Variable}: Attack\_type (categorical)

  \begin{itemize}
  \tightlist
  \item
    Attack patterns:

    \begin{itemize}
    \tightlist
    \item
      DOS\_SYN\_Hping
    \item
      ARP\_poisioning
    \item
      NMAP\_UDP\_SCAN
    \item
      NMAP\_XMAS\_TREE\_SCAN
    \item
      NMAP\_OS\_DETECTION
    \item
      NMAP\_TCP\_scan
    \item
      DDOS\_Slowloris
    \item
      Metasploit\_Brute\_Force\_SSH
    \item
      NMAP\_FIN\_SCAN
    \end{itemize}
  \item
    Normal Patterns:

    \begin{itemize}
    \tightlist
    \item
      MQTT
    \item
      Thing\_speak
    \item
      Wipro\_bulb\_Dataset
    \item
      Amazon-Alexa (\textbf{MISSING FROM DATASET})
    \end{itemize}
  \end{itemize}
\end{enumerate}

\section{Previous Work:}\label{previous-work}

\section{Report Organization:}\label{report-organization}

The remainder of this report is structured as follows.

Section 2 presents an exploratory analysis of the RT-IoT 2022 dataset,
examining feature distributions, statistical separation between normal
and attack traffic, and the overall structure of the data through
techniques such as kernel density estimation, skewness analysis,
correlation mapping, and principal component analysis. Diagnostic tools
such as variance inflation factors are also used to examine
multicollinearity among predictors. These analyses provide insight into
the behavioral differences between benign and malicious network flows
and motivate the preprocessing and modeling choices used in later
sections.

Section 3 describes the feature preprocessing pipeline used to prepare
the dataset for modeling. This includes the encoding of categorical
variables, identification of highly skewed features, and the application
of log transformations to stabilize feature distributions.

Section 4 introduces a supervised baseline classifier using Logistic
Regression. This section outlines the modeling framework, training
procedure, and evaluation metrics used to establish a statistical
benchmark for intrusion detection performance on the dataset.

Section 5 extends the analysis to unsupervised anomaly detection using
autoencoder-based models. These models are trained on normal traffic
patterns and evaluated on their ability to identify anomalous behavior
corresponding to attack traffic.

Section 6 compares the performance of the supervised and unsupervised
approaches and discusses their strengths and limitations.

Finally, Section 7 concludes the report by summarizing the key findings,
discussing dataset limitations---including the absence of Amazon Alexa
traffic---and outlining directions for future work, including edge
deployment considerations for IoT intrusion detection systems.

\bookmarksetup{startatroot}

\chapter{Exploratory Data Analysis
(EDA):}\label{exploratory-data-analysis-eda}

\begin{center}\rule{0.5\linewidth}{0.5pt}\end{center}

\bookmarksetup{startatroot}

\chapter{Preprocessing}\label{preprocessing}

\begin{center}\rule{0.5\linewidth}{0.5pt}\end{center}

\bookmarksetup{startatroot}

\chapter{Baseline Logistic Regression
Model}\label{baseline-logistic-regression-model}

\begin{center}\rule{0.5\linewidth}{0.5pt}\end{center}

\bookmarksetup{startatroot}

\chapter{Autoencoder Models}\label{autoencoder-models}

\begin{center}\rule{0.5\linewidth}{0.5pt}\end{center}

\bookmarksetup{startatroot}

\chapter{Model Comparison}\label{model-comparison}

\begin{center}\rule{0.5\linewidth}{0.5pt}\end{center}

\bookmarksetup{startatroot}

\chapter{Conclusion}\label{conclusion}




\end{document}
